
\textbf{Exemple 1.5.} Soit $S$ un schéma localement noethérien universellement caténaire (par exemple sous-schéma d'un schéma de Cohen-Macaulay). Posons
\begin{equation}\tag{1.5.1}
D_S(x)=n-\codim(\overline{\{x\}},S)=n-\dimdot\underline{O}_{S,x}\qquad(x\in S),
\end{equation}
et on pose aussi pour tout $S$-schéma propre $s_Y:Y\to S$ et $y\in Y$ :
\begin{equation}\tag{1.5.2}
D_Y(y)=D_S(s_Y(y))+\degtr_{k(s_Y(y))}k(y).
\end{equation}

Si $f:Y_1\to Y_2$ est un $S$-morphisme de $S$-schémas propres, l'additivité des degrés de transcendance et (1.5.2) donnent immédiatement
\begin{equation}\tag{1.5.3}
D_{Y_1}(y)=D_{Y_2}(f(y))+\degtr_{k(f(y))}k(y),\qquad(y\in Y_1).
\end{equation}
Montrons que les fonctions $D_Y$ ainsi définies satisfont aux conditions (i)-(iv) de 1.4. D'après (1.5.3), la condition (ii) est la seule non-triviale. Grâce à (iv) on peut évidemment supposer que $Y=\overline{\{x\}}$. On factorise alors le morphisme structural $s_Y$ de $Y$ à travers son image $s_Y(Y)=Z\hookrightarrow S$, soit $s_Y=g\circ f$ :
\[
\begin{tikzcd}
Y \arrow[r,"f"] & Z \arrow[r,hook,"g"] & S,
\end{tikzcd}
\]
où $Y$ et $Z$ sont irréductibles, $f$ propre et surjectif et $g$ une immersion fermée. Posons $x'=s_Y(x)$, $y'=s_Y(y)$. On déduit
alors de (1.5.2) et (1.5.1) :
\[
D_Y(x)-D_Y(y)=\dimdot\underline{O}_{S,y'}-\dimdot\underline{O}_{S,x'}+\degtr_{k(x')}k(x)-\degtr_{k(y')}k(y).
\]
Puisque $S$ est caténaire, on a
